% format de la date pour Git
\DTMnewdatestyle{mydate}{%
  \renewcommand*{\DTMdisplaydate}[4]{##3 \DTMfrenchMonthname{\DTMtwodigits{##2}} ##1}%
  \renewcommand{\DTMDisplaydate}{\DTMdisplaydate}%
}
\DTMsetdatestyle{mydate}

% definition du format de l'entete et du pied de page
\pagestyle{fancyplain}
\lhead{}
\rhead{\leftmark}
\rfoot{
  \begin{scriptsize}
  \textsl{version du \DTMusedate{gitdate}}
  \end{scriptsize}
}

%%% warning encadré
%%% utilisation : \begin{warning} blablabla \end{warning}
\usepackage{pifont,mdframed}
\newenvironment{warning}
  %%{\par\begin{mdframed}[linewidth=2pt,linecolor=red]%
  {\par\begin{mdframed}[linewidth=2pt]%
    \begin{list}{}{\leftmargin=1cm
                   \labelwidth=\leftmargin}\item[\Large\ding{43}]}
  {\end{list}\end{mdframed}\par}

%%% pour l'usage des programmes qui ne tiennent pas sur une ligne
%%% utilisation :
%%%%%% \minibox[frame]{
%%%%%%  \texttt{gengraphvect\.py [-h] -a ACQID [-o FORMAT] [-n NIVEAU]} \\
%%%%%%  \rightline{ \texttt{[-f FICHIER] [-s SENS] [-v] [VECTEUR ...]} }
%%%%%% }
\usepackage{minibox}
